\section{Benchmarks for Klauer's task examples}
\label{sec:annex-p19-rationality-2-5}

\begin{annexitem}{annex:p19.r25.benchmarks}{thm:p19-r25-benchmarks}{Rule testing, stipulated deduction, and posterior odds}
\begin{proof}
Let \(p\) mean that a card's letter is \(A\) and \(q\) that its
digit is \(3\). The classical \(p\to q\) rule is falsified exactly by a card whose
letter is \(A\) and whose digit is not \(3\). The visible \(A\)
could hide a non-\(3\) digit and must be inspected. The visible
\(1\) could hide an \(A\) and must be inspected. A visible
\(B\) has a non-\(A\) letter and cannot be such a falsifier.
A visible \(3\) satisfies the consequent regardless of its hidden
letter and cannot be a falsifier. Thus the two specified cards, and
only those two, test every displayed possibility of a violation.
Selecting the two symbols named by the rule is a different strategy.

If \(c\) means that the child is crying and \(h\) that the child
is happy, then any valuation satisfying \(c\to h\) and \(c\)
satisfies \(h\). A valuation with \(c\) true and \(h\) false
would falsify the supplied conditional, so is excluded by the task.
Likewise the finishing/fullness premises entail fullness; under the
training example's stipulated incompatible full/empty alternatives,
they exclude emptiness. These are semantic benchmark checks, not
evidence that a participant enumerated valuations or used a particular
cognitive algorithm.

For the two exhaustive categories let their priors be \(5/1004\)
and \(999/1004\), and write \(\ell_P,\ell_N>0\) for the description
likelihoods. Bayes' rule gives posterior weights proportional to
\(5\ell_P\) and \(999\ell_N\). Dividing gives the displayed
odds \(5\ell_P/(999\ell_N)\), which exceed one exactly when
\(\ell_P/\ell_N>999/5\). Choosing likelihood ratio one leaves
the nanny category dominant; choosing ratio 300 makes the politician
category dominant. Thus the given category counts do not alone settle
the full conditional judgment. The assumed likelihoods are stipulated
comparison models, not measured stereotype strengths in the cited
experiment. This benchmark bridge uses the finite Bayesian identity
at \cref{thm:p15-finite-bayes}; it does not attribute a new experiment
or posterior calculation to Klauer's figure.
\end{proof}
\end{annexitem}
